\documentclass[main.tex]{subfiles}


\begin{document}

%from pagenumber to up
% \phantomsection 
% \hypertarget{toc}{}

 

 

% номер секции вручную
\setcounter{section}{29
}
\addtocounter{section}{-1} 

\section{Закон Архимеда}


\sbox{0}{%
    \begin{tikzpicture}[scale=1,font=\small,            line cap=round,                             >={Stealth[inset=0pt,round,length=3mm, angle'=20]},vec/.style={>={Stealth[inset=0pt,round,length=3.5mm, angle'=20]}}]


%liquid
\fill[SkyBlue,semitransparent] (0,0) --++ (0,2.5) --++ (2,0) --++ (0,-.7) --++ (1,0) --++ (0,.7) --++ (2,0) --++ (0,-2.5) -- cycle;

\begin{scope}[yshift=.3cm]
    
\fill[yellow!70!black,semitransparent] (2,2.5) rectangle++ (1,-1);
\draw[] (2,2.5) rectangle++ (1,-1);

%forces
\draw[->,red,thick,vec] (2.5,2) --++ (0,1) node[black,right] {$\vec F_\text{А}$};
\node at (2.5,2) [circle,fill=red,inner sep=0pt,minimum size=1mm,label=] {};

\draw[->,red,thick,vec] (2.5,2) --++ (0,-1) node[black,right] {$\vec F_\text{т}$};

\end{scope}


    \end{tikzpicture}%
}

{%
\setlength\intextsep{0pt}
\begin{wrapfigure}[10]{r}{\wd0}
    \centering
    \usebox{0}%
    \caption{Куб на плаву}
    \label{fig:p53}
\end{wrapfigure}


Некоторые тела, помещенные в жидкость\footnote{Далее считается, что среда (жидкость или газ) покоится у поверхности планеты.}, не тонут. В таких случаях сила тяжести уравновешивается какой-то другой силой, действующей на тело со стороны жидкости. Эта сила называется \emph{выталкивающей} или \emph{архимедовой} силой и действует на любое тело, погруженное в жидкость или газ целиком или частично.

На рис.\,\ref{fig:p53} изображен деревянный куб, покоящийся на поверхности воды; сила тяжести равна так называемой силе~Архимеда: $F_\text{т}=F_\text{А}$. 

}

\begin{law}
\textbf{Закон Архимеда.} На погружённое (возможно, частично) в жидкость или газ тело действует выталкивающая сила, равная весу среды, объём которой вытеснило тело:
\begin{equation}
    F_\text{А}=P_\text{выт.\,с},
\end{equation}
где~$P_\text{выт.\,с}$~--- вес вытесненной среды. (Силы связаны третьим законом Ньютона.)
\end{law}
%
Можно заметить, что вес вытесненной среды в рассматриваемых условиях есть $P_\text{выт.\,с}=\rho_\text{с}gV_\text{выт}$; где~$\rho_\text{с}$~--- плотность среды, $V_\text{выт}$~--- вытесненный объем. Тогда сила Архимеда равна:
\begin{equation}
    F_\text{А}=\rho_\text{с}gV_\text{выт}.
\end{equation}

\emph{Плавание}~--- это состояние тела, при котором оно не тонет в жидкости (или газе), будучи погруженным в нее. На рис.\,\ref{fig:p54} показаны три погруженных в воду шара одинакового размера, сделанных из разных материалов.

\begin{figure}[H]
    \centering

    \begin{tikzpicture}[font=\small,            line cap=round,                             >={Stealth[inset=0pt,round,length=3mm, angle'=20]},vec/.style={>={Stealth[inset=0pt,round,length=3.5mm, angle'=20]}}]


\fill[SkyBlue,semitransparent] (0,-.27) rectangle (12,2.5);

%balls
\shade[ball color=lightgray] (3,1.25) node[xshift=-.3cm,left] {Ш$_1$} circle (.3);
\shade[ball color=SkyBlue!70!RoyalBlue] (6,1.25) node[xshift=-.3cm,left] {Ш$_2$} circle (.3);
\shade[ball color=yellow!40!brown] (9,1.25) node[xshift=-.3cm,left] {Ш$_3$} circle (.3);

%F
\def\w{.9}
%легкий бетон
\draw[->,red,thick,vec] (3,1.25) --++ (0,-1.2) node[black,right] {$\vec F_\text{т1}$};
\draw[->,red,thick,vec] (3,1.25) --++ (0,\w) node[black,right] {$\vec F_\text{А}$};
\node at (3,1.25) [circle,fill=red,inner sep=0pt,minimum size=1mm,label=] {};

%water
\draw[->,red,thick,vec] (6,1.25) --++ (0,-\w) node[black,right] {$\vec F_\text{т2}$};
\draw[->,red,thick,vec] (6,1.25) --++ (0,\w) node[black,right] {$\vec F_\text{А}$};
\node at (6,1.25) [circle,fill=red,inner sep=0pt,minimum size=1mm,label=] {};

%дерево
\draw[->,red,thick,vec] (9,1.25) --++ (0,-.7) node[black,right] {$\vec F_\text{т3}$};
\draw[->,red,thick,vec] (9,1.25) --++ (0,\w) node[black,right] {$\vec F_\text{А}$};
\node at (9,1.25) [circle,fill=red,inner sep=0pt,minimum size=1mm,label=] {};


    \end{tikzpicture}
\caption{Шары в воде}
\label{fig:p54}
\end{figure}

На примере с шарами~Ш$_1$, Ш$_2$ и~Ш$_3$  плотностей~$\rho_1$, $\rho_2$ и~$\rho_3$ соответственно, которые вначале покоятся в жидкости плотности~$\rho_\text{с}$, можно проиллюстрировать три возможных движения тела после погружения в некоторую среду.

\begin{enumerate}
    \item Шар~Ш$_1$ сделан из бетона: $F_\text{т1}>F_\text{А}$ или $\rho_1>\rho_\text{с}$. Этот шар тонет.
    \item В шаре~Ш$_2$ вода, обурнутая легкой тонкой пленкой: $F_\text{т2}=F_\text{А}$ или $\rho_2=\rho_\text{с}$. Этот шар остается в покое.
    \item Шар~Ш$_3$ деревянный: $F_\text{т3}<F_\text{А}$ или $\rho_3<\rho_\text{с}$. Этот шар всплывает. Он придет в равновесие у поверхности жидкости, частично погрузившись в нее.
    \looseness=-1
\end{enumerate}


\emph{Условие плавания} тела можно записать в виде неравенства: $\rho \leqslant \rho_\text{с}$, где~$\rho$~--- плотность тела.








 
\end{document}